\chapter{Comónadas y Coálgebras}

En este capítulo daremos una introducción a los conceptos de coálgebra y comónada. El resultado
principal, y más relevante para la presente tesis, de esta sección es que, bajo ciertas condiciones, la categoría de coálgebras para comónada es un topos.

\section{Comónadas}

Las comónadas, como sugiere la terminología, son estructuras duales a las mónadas. 

\begin{definition}
  Una comónada en una categoría $\catC$ es un triplete $(F,\epsilon,\delta)$ donde 
  $F\colon \catC\to \catC$ es un funtor y $\epsilon\colon F\to \id_{\catC}$ 
  y $\delta\colon F\to F^2$ son transformaciones naturales, llamadas counidad y comultiplicación respectivamente, tales que los siguientes diagramas conmutan:
  \begin{equation*}
    \begin{tikzcd}
      {} & F\ar[swap]{ld}{\Id}\ar{d}{\delta}\ar{rd}{\Id} & {} & {} & F\ar{r}{\delta}\ar[swap]{d}{\delta} & F^2\ar{d}{\delta_F} \\
      F & F^2\ar{l}{\epsilon_F}\ar[swap]{r}{F(\epsilon)} & F, & {} & F^2\ar[swap]{r}{F(\delta)} & F^3.
    \end{tikzcd}
  \end{equation*}
\end{definition}

En los diagramas duales de mónada se puede apreciar el parentesco con los axiomas de monoide.

\begin{proposition}[label=prop:comonadainducida]
  Dada una adjunción $L\dashv R$, donde $L\colon\catC\to\catD$ y $R\colon\catD\to\catC$, con unidad $\eta\colon \id_{\catC}\to RL$ y counidad 
  $\epsilon\colon LR\to \id_{\catD}$, entonces $(LR,\epsilon,L(\eta_R))$ es una comónada en $\catD$.
\end{proposition}
\begin{proof}
  Primero probemos la counidad. Queremos mostrar la conmutatividad de ambos triángulos en el diagrama:
  \begin{equation*}
    \begin{tikzcd}
      {} & LR\ar[swap]{ld}{\Id}\ar{d}{L(\eta_R)}\ar{rd}{\Id} & {} \\
      LR & LRLR\ar{l}{\epsilon_{LR}}\ar[swap]{r}{LR(\epsilon)} & LR.
    \end{tikzcd}
  \end{equation*}
  El lado izquierdo es la identidad triangular $\epsilon_L\circ L(\eta) = 1_L$ en la imagen de $R$.
  Similarmente, el lado derecho resulta de aplicar $L$ a la identidad triangular $R(\epsilon)\circ\eta_R = 1_R$.
  Ahora, para demostrar la coasociatividad, tomamos $X$ en $\catC$ y consideramos el siguiente diagrama:
  \begin{equation*}
    \begin{tikzcd}
      X\ar[swap]{d}{\eta_X}\ar{r}{\eta_X} & RLX\ar{d}{\eta_{RLX}} \\
      RLX\ar[swap]{r}{RL(\eta_X)} & RLRLX.
    \end{tikzcd}
  \end{equation*}
  Que es un cuadro de naturalidad de $\eta$ aplicada al morfismo $\eta_X$. Notemos que si
  introducimos $R$ por la derecha y aplicamos $L$ al diagrama, obtenemos al diagrama de coasociatividad.
\end{proof}

\section{Coálgebras para una comónada}

\begin{definition}
  Dada una categoría $\catC$ y un endofuntor $F$, una $F$-Coálgebra es un par $(A,s)$ donde $A$ es un objeto de $\catC$ y $s\colon A\to FA$ un morfismo, a menudo llamado el "morfismo estructural".
\end{definition}

\begin{definition}
  Un morfismo de coálgebras entre $(A,s)$ y $(B,t)$ es un morfismo $f\colon A\to B$ en $\catC$ tal que el siguiente diagrama conmuta:
  \begin{equation*}
    \begin{tikzcd}
      A\ar[swap]{d}{s}\ar{r}{f} & B\ar{d}{t} \\
      FA\ar[swap]{r}{F(f)} & FB.
    \end{tikzcd}
  \end{equation*}
\end{definition}

\begin{proposition}
  Dada una categoría $\catC$ y un endofuntor $F$, las $F$-coálgebras y sus morfismos forman una categoría.
\end{proposition}
\begin{proof}
  Dada cualquier coálgebra $(A,s)$, la identidad $\id_A$ es un morfismo de coálgebras, ya que
  $F(\id_A) = \id_{FA}$. Para componer morfismos pegamos los cuadrados conmutativos correspondientes 
  y de nuevo la funtorialidad de $F$ nos asegura el resultado. La asociatividad de la composición se 
  hereda directamente de $\catC$.
\end{proof}

Análogamente a como se construye la categoría de álgebras para una mónada para demostrar que toda mónada viene de una adjunción,
se puede realizar una construcción dual para comónadas.

\begin{definition}
  Dada una comónada $\textbf{F}=(F,\epsilon,\delta)$ en una categoría $\catC$, la categoría de coálgebras para la comónada $\catC_{\textbf{F}}$ 
  está constituida por los morfismos entre $F$-coálgebras $(A,s)$ tal que satisfacen las condiciones:
  \begin{equation*}
    \begin{tikzcd}
      A & FA\ar[swap]{l}{\epsilon_A} & {} & F^2A & FA\ar[swap]{l}{F(s)} \\
      {} & A,\ar{ul}{\id_A}\ar[swap]{u}{s} & {} & FA\ar{u}{\delta_A} & A.\ar{l}{s}\ar[swap]{u}{s}
    \end{tikzcd}
  \end{equation*}
\end{definition}

\begin{proposition}[label=prop:adjcolibreolvida]
  Si $\textbf{F}=(F,\epsilon,\delta)$ es una comónada en $\catC$, entonces hay una adjunción 
  $U\dashv R$ donde $U\colon \catC_{\textbf{F}}\to \catC$ es el funtor olvidadizo y $R\colon \catC\to \catC_{\textbf{F}}$ 
  es el funtor "colibre".
\end{proposition}
\begin{proof}
  En objetos el funtor $U$ se olvida de la estructura de coálgebra y en morfismos olvida que
  se trataba de un morfismo entre coálgebras. Por su parte, definimos a $R$ en objetos $A\in\catC$
  como $RA = (FA,\delta_A\colon FA\to F^2A)$, aprovechando la estructura de comónada. Entonces en morfismos
  $g\colon A\to B$ proponemos $R(g) = F(g)$ puesto que el diagrama correspondiente
  \begin{equation*}
    \begin{tikzcd}
      FA\ar[swap]{d}{\delta_A}\ar{r}{F(g)} & FB\ar{d}{\delta_B} \\
      F^2A\ar[swap]{r}{F^2(g)} & F^B,
    \end{tikzcd}
  \end{equation*}
  es un cuadro de naturalidad de $\delta$, así que $F(g)$ es un morfismo de coálgebras y
  se puede comprobar la funtorialidad de la asignación usando que $F$ es funtor. Más aun,
  la asignación es colibre en el sentido de que, para toda $(X,t)$ coálgebra y $f\colon X\to A$
  morfismo, existe un único morfismo $\overline{f}\colon(X,t)\to(FA,\delta_A)$ tal que $\epsilon_AU(\overline{f})=f$.
  Primero, proponemos $\overline{f}=F(f)t$, que hace conmutar al diagrama:
  \begin{equation*}
    \begin{tikzcd}
      X\ar[swap]{d}{t}\ar{r}{F(f)t} & FA\ar{d}{\delta_A} \\
      FX\ar[swap]{r}{F^2(f)F(t)} & F^2A.
    \end{tikzcd}
  \end{equation*}
  Para demostrar la conmutatividad usamos que $(X,t)$ es coálgebra para la comónada entonces, $F^2(f)F(t)t=F^2(f)\delta_Xt$ y por la naturalidad
  de $\delta$, $F^2(f)\delta_Xt=\delta_AF(f)t$, demostrando que $F(f)t$ es morfismo
  de coálgebras. Luego, si hubiera otro $g\colon(X,t)\to(FA,\delta_A)$ que cumpla la propiedad,
  pasaría que el diagrama
  \begin{equation*}
    \begin{tikzcd}
      X\ar[shift left = 0.4ex]{d}{\overline{f}}\ar[swap,shift right = 0.4ex]{d}{g}\ar{r}{t} & FX\ar[shift left = 0.4ex]{d}{F\overline{f}}\ar[swap,shift right = 0.4ex]{d}{Fg}\ar{rd}{F(f)} & {}\\
      FA\ar[swap]{r}{\delta_A} & F^2A\ar[swap]{r}{F(\epsilon_A)} & FA,
    \end{tikzcd}
  \end{equation*}
  conmuta. Pero $F(\epsilon_A)\delta_A=\id$ por la condición de counidad, entonces $\overline{f} = F(f)t = g$, demostrando la unicidad.
  Reescribiendo, lo que demostramos es que para todo $A\in\catC$ y $X\in\catC_\textbf{F}$ con $f\colon UX\to A$, el siguiente diagrama conmuta:
  \begin{equation*}
    \begin{tikzcd}
      UX\ar[swap]{rd}{f}\ar[dashed]{r}{\overline{f}} & URA\ar{d}{\epsilon_A} \\
      & A.
    \end{tikzcd}
  \end{equation*}
  Es decir, $\epsilon$ es couniversal y tenemos que $U\dashv R$.
\end{proof}

\begin{theorem}
  Dada una adjunción $L\colon\catC\leftrightarrows\catD :R$, si $G=L\circ R$, entonces existe un funtor de "comparación"
  $K\colon\catC\to\catD_{\textbf{G}}$ tal que $U\circ K=L$ y $K\circ R=F$.
\end{theorem}
\begin{proof}
  Llamemos $\eta$ y $\epsilon$ a la unidad y counidad respectivamente, de la adjunción dada $L\dashv R$, 
  entonces ya sabemos que $(G,\epsilon,L(\eta_R))$ es una comónada en $\catD$ por \ref{prop:comonadainducida} y podemos construir a $\catD_{\textbf{G}}$.
  Si $F$ es el funtor colibre y $U$ el que olvida, tenemos una situación así:
  \begin{equation*}
    \begin{tikzcd}
      \catC\ar[shift left = 0.4ex]{d}{L} & \catD_{\textbf{G}}\ar[shift left = 0.4ex]{d}{U} \\
      \catD\ar[shift left = 0.4ex]{u}{R}\ar[no head,shift right = 0.4ex]{r}{}\ar[no head,shift left = 0.4ex]{r}{} & \catD.\ar[shift left = 0.4ex]{u}{F}
    \end{tikzcd}
  \end{equation*}
  De manera dual a la situación con mónadas, es posible definir un funtor $K\colon\catC\to\catD_{\textbf{G}}$
  que conmute tanto con los adjuntos izquierdos como con los derechos y tal que exhibe una equivalencia de categorías.
  Proponemos $K(A) = (LA,L\eta_A)$ y si $f\colon A\to B$, $K(f) = L(f)$, pero primero 
  necesitamos que $(LA,L\eta_A)$ sea una coálgebra para la comónada, para ello consideramos a los diagramas:
  \begin{equation*}
    \begin{tikzcd}
      LA\ar{r}{L\eta_A}\ar[swap]{rd}{\id} & LRLA\ar{d}{\epsilon_{LA}} & {} & LA\ar{r}{L\eta_A}\ar[swap]{d}{L\eta_A} & GLA\ar{d}{L\eta_{RA}} \\
      & LA, & {} & GLA\ar[swap]{r}{GL\eta_A} & G^2LA.
    \end{tikzcd}
  \end{equation*}
  El diagrama izquierdo conmuta al ser una de las identidades triangulares y el diagrama 
  derecho es un cuadro de naturalidad de $L\eta$, entonces la coálgebra $K(A)$ es compatible 
  con la comónada. Por otro lado, para verificar que $K(f)$ es un morfismo de coálgebras, el 
  siguiente diagrama debe conmutar:
  \begin{equation*}
    \begin{tikzcd}
      LA\ar{r}{Lf}\ar[swap]{d}{L\alpha_A} & LB\ar{d}{L\alpha_B} \\
      LRLA\ar[swap]{r}{L^2f} & LRLB.
    \end{tikzcd}
  \end{equation*}
  Pero nuevamente es un cuadro de conmutatividad de $L\eta$. Luego, por como definimos a $K$,
  es inmediato comprobar que $UK=L$ y $KR=F$. 
\end{proof}

\begin{theorem}
  Sean $\mathcal{E}$ un topos y $\textbf{G}=(G,\epsilon,\delta)$ una comonada en $\mathcal{E}$ tal que
  $G$ es exacto por la izquierda. Entonces la categoría de coálgebras $\mathcal{E}_\textbf{G}$ es un topos.
\end{theorem}
\begin{lemma}
  La categoría de coálgebras $\mathcal{E}_\textbf{G}$ tiene límites finitos.
\end{lemma}
\begin{proof}
  Es suficiente mostrar que la categoría de coálgebras tiene productos e igualadores finitos. 
  Como $\mathcal{E}$ tiene productos y $G$ es funtor, dadas $(A,s),(B,t)$ coálgebras de $C$ tenemos:
  \begin{equation*}
    \begin{tikzcd}
        GA & GA\times GB \ar[swap]{l}{\pi_{GA}} \ar{r}{\pi_{GB}} & GB\\
        & G(A\times B). \ar{ul}{G(\pi_A)} \ar[dashed]{u}{\phi} \ar[swap]{ur}{G(\pi_B)}
    \end{tikzcd}
  \end{equation*}
  Con $\phi$ natural en los objetos A y B. En este caso, gracias a que $G$ es
  exacto izquierdo, $\phi$ es un isomorfismo. Proponemos
  $(A\times B, \phi^{-1}(s\times t))$ como el producto de $(A,s)$ y $(B,t)$. Veamos que
  es una coálgebra de $C$.
  \begin{equation*}
    \begin{tikzcd}
        A\times B\ar{r}{s\times t}\ar[swap]{d}{\pi_A}\ar[phantom]{2-2}{(a)} & GA\times GB \ar[swap]{d}{\pi_{GA}}\ar{r}{\phi^{-1}}\ar[phantom]{2-3}[near start]{(b)} 
          & G(A\times B)\ar{dl}{G(\pi_A)}\ar{d}{\epsilon_{A\times B}} \\
        A\ar{r}{s}\ar[swap]{dr}{\id_A} & GA\ar{d}{\epsilon_A}\ar[phantom]{r}{(d)}\ar[phantom]{3-1}[near start]{(c)} & A\times B \ar{dl}{\pi_A} \\
        {} & A.
    \end{tikzcd}
  \end{equation*}
  En el diagrama (a) conmuta por la definición de $s\times t$, (b) conmuta gracias a la construcción
  de $\phi$, (c) conmuta ya que $(A,s)$ es coálgebra y (d) conmuta pues es un cuadrado de naturalidad
  de $\epsilon$. Por tanto $\pi_A=\pi_A \epsilon_{A\times B} \phi^{-1}(s\times t)$ y dado que podemos
  repetir el argumento para $\pi_B$, concluimos que $\epsilon_{A\times B} \phi^{-1}(s\times t)$
  es la identidad en $A\times B$. Esto es la ecuación de la counidad, entonces seguimos con la
  coasociatividad. De manera similar a $\phi$ podemos encontrar un isomorfismo $\psi\colon G(GA\times GB) \to G^2A\times G^2B$,
  lo cual nos da el siguiente diagrama de isomorfismos,
  \begin{equation*}
    \begin{tikzcd}
      G(GA\times GB)\ar{r}{G(\phi^{-1})} & G^2(A\times B) \\
      G^2A\times G^2B.\ar{u}{\psi^{-1}}\ar[swap]{ur}{G(\phi^{-1})\psi^{-1}}
    \end{tikzcd}
  \end{equation*}
  Para la coasociatividad queremos ver la conmutatividad del siguiente diagrama:
  \begin{equation*}
    \begin{tikzcd}
      A\times B\ar[phantom]{rd}{(a)}\ar{r}{s\times t}\ar[swap]{d}{s\times t} & GA\times GB\ar{r}{\phi^{-1}}\ar{d}{\delta_A\times \delta_B}
        & G(A\times B)\ar{dd}{\delta_{A\times B}} \\
      GA\times GB\ar{r}{G(s)\times G(t)}\ar[swap]{d}{\phi^{-1}} & G^2A\times G^2B\ar{d}{\psi^{-1}} \\
      G(A\times B)\ar[phantom]{ur}{(b)}\ar[swap]{r}{G(s\times t)} & G(GA\times GB)\ar[phantom]{uur}{(c)}\ar{r}{G(\phi^{-1})} & G^2(A\times B).
    \end{tikzcd}
  \end{equation*}
  En (a) sabemos que el diagrama en cada coordenada conmuta pues $(A,s),(B,t)$ son coálgebras. 
  Para (b) tenemos:
  \begin{equation*}
    \begin{tikzcd}
      GA\ar[swap]{d}{G(s)} & GA\times GB\ar[swap]{l}{\pi_{GA}}\ar{r}{\pi_{GB}}\ar[dashed]{d}{G(s)\times G(t)} & GB\ar{d}{G(t)}\\
      G^2A & G^2A\times G^2B\ar[swap]{l}{\pi_{G^2A}}\ar{r}{\pi_{G^2B}} & G^2B \\
      G^2A\ar{u}{\id_{G^2A}} & G(GA\times GB)\ar[swap]{l}{G(\pi_{GA})}\ar{u}{\psi}\ar{r}{G(\pi_{GB})} & G^2B\ar[swap]{u}{\id_{G^2B}} \\
      GA\ar{u}{G(s)} & G(A\times B)\ar[swap]{l}{G(\pi_A)}\ar{u}{G(s\times t)}\ar{r}{G(\pi_B)} & GB\ar[swap]{u}{G(t)} \\
      & GA\times GB.\ar{ul}{\pi_{GA}}\ar{u}{\phi^{-1}}\ar[swap]{ur}{\pi_{GB}}
    \end{tikzcd}
  \end{equation*}
  Cada parte conmuta gracias a cómo construimos los morfismos usando la propiedad universal del producto,
  por lo que $G(s)\times G(t)=\langle G(s)\pi_{GA},G(t)\pi_{GB} \rangle = \psi G(s\times t)\phi^{-1}$,
  lo cual nos da la identidad deseada al poscomponer con $\psi^{-1}$. En el caso de (c) observamos que:
  \begin{equation*}
    \begin{tikzcd}
      {} & GB\ar{r}{\delta_B} & G^2B & {} \\
      GA\times GB\ar{ur}{\pi_{GB}}\ar{r}{\phi^{-1}}\ar[swap]{dr}{\pi_A} & G(A\times B)\ar[swap]{u}{G(\pi_B)}\ar{r}{\delta_{A\times B}}\ar{d}{G(\pi_A)}
        & G^2(A\times B)\ar{u}{G^2(\pi_B)}\ar{r}{\psi G(\phi)}\ar[swap]{d}{G^2(\pi_A)}
        & G^2A\times G^2GB\ar[swap]{ul}{\pi_{G^2B}}\ar{dl}{\pi_{G^2A}} \\
      & GA\ar[swap]{r}{\delta_A} & G^2A,
    \end{tikzcd}
  \end{equation*}
  donde todo conmuta ya que son diagramas de productos y cuadrados de naturalidad de $\delta$, así
  que $\delta_A\times\delta_B=\psi G(\phi)\delta_{A\times B}\phi^{-1}$, pero $\psi G(\phi)$ es un
  isomorfismo, entonces obtenemos la igualdad que queríamos.

  Ahora consideramos a los morfismos de coálgebras $f\colon (X,u)\to (A,s)$ y $g\colon (X,u)\to (B,t)$ para demostrar
  la propiedad universal de $(A\times B, \phi^{-1}(s\times t))$. Notemos que podemos construir el siguiente diagrama:
  \begin{equation*}
    \begin{tikzcd}
      X\ar[swap]{d}{u}\ar{r}{f\times g} & A\times B\ar{dr}{\phi^{-1}(s\times t)}\ar{d}{s\times t} \\
      GX\ar[swap]{r}{G(f)\times G(g)} & GA\times GB\ar[swap]{r}{\phi^{-1}} & G(A\times B).
    \end{tikzcd}
  \end{equation*}
  El cuadro conmuta pues $f$ y $g$ son morfismos de coálgebras, además el diagrama exterior sería un 
  diagrama de morfismo de coálgebras si tuviéramos que $\phi^{-1}(G(f)\times G(g))=G(f\times g)$,
  pero, usando la propiedad universal de $G(A\times B)$ como producto, es sencillo comprobar dicha igualdad:
  \begin{equation*}
    \begin{tikzcd}
      {} & {GA}\ar{r}{\text{id}} & GA \\
      GX\ar[bend left = 5ex]{ru}{Gf}\ar[swap, bend right = 5ex]{rd}{Gg}\ar{r}{Gf\times Gg} & GA\times GB\ar{r}{\phi^{-1}}\ar{u}{\pi_{GA}}\ar[swap]{d}{\pi_{GB}} & G(A\times B)\ar{u}{G(\pi_A)}\ar[swap]{d}{G(\pi_B)} & GX\ar[swap, bend right = 5ex]{lu}{Gf}\ar[bend left = 5ex]{ld}{Gg}\ar{l}{G(f\times g)}\\
      {} & {GB}\ar[swap]{r}{\text{id}} & GB \\
    \end{tikzcd}
  \end{equation*}
  Cada parte de la izquierda conmuta por definición de $\phi^{-1}$ y la de $G(f)\times G(g)$,
  y la derecha conmuta al ser $G$ aplicado a la construcción de $f\times g$, así que, por la propiedad universal del producto, concluimos que $\phi^{-1}(G(f)\times G(g))=G(f\times g)$. Con esto logramos
  exhibir un morfismo $f\times g\colon (X,u)\to (A\times B,\phi^{-1}(s\times t))$. La conmutatividad con las proyecciones es inmediataLa y la unicidad
  se da gracias a la unicidad de $f\times g$ en $\mathcal{E}$. Por tanto $\mathcal{E}_\textbf{G}$
  tiene productos finitos y procedemos a mostrar que tiene igualadores. Consideramos $(A,s)$,
  $(B,t)$ coálgebras y $f,g\colon (A,s)\to (B,t)$. Proponemos $(E,u)$ como sigue:
  \begin{equation*}
    \begin{tikzcd}
      E\ar{r}{e}\ar[swap,dashed]{d}{u} & A\ar[shift left=0.75ex]{r}{f}\ar[swap, shift right=0.75ex]{r}{g}\ar{d}{s} & B\ar{d}{t} \\
      GE\ar[swap]{r}{G(e)} & GA\ar[shift left=0.75ex]{r}{G(f)}\ar[swap, shift right=0.75ex]{r}{G(g)} & GB. \\
    \end{tikzcd}
  \end{equation*}
  Los cuadros conmutan ya que $f$ y $g$ son morfismos de coálgebras, entonces $s\circ e$ iguala
  a $G(f)$ con $G(g)$, lo cual nos da la existencia de $u$. $GE$ es el igualador de $G(f)$ y $G(g)$ salvo isomorfismo
  porque $G$ preserva límites. Veamos que $(E,u)$ es coálgebra, para ello empezamos con la ``counidad'',
  es decir, queremos demostrar la conmutatividad del siguiente diagrama:
  \begin{equation*}
    \begin{tikzcd}
      E\ar{r}{u}\ar[swap]{dr}{\id_E} & GE\ar{d}{\epsilon_E} \\
      & E
    \end{tikzcd}
  \end{equation*}
  Notemos que $\epsilon_AG(e)$ iguala a $f,g$ por la naturalidad de $\epsilon$,
  como se muestra en el diagrama de abajo a la izquierda.
  \begin{equation*}
    \begin{tikzcd}
      E\ar{r}{e} & A\ar[shift left=0.75ex]{r}{f}\ar[swap, shift right=0.75ex]{r}{g} & B & {} E\ar{r}{e}\ar[swap]{d}{u} & A\ar{r}{\id_A}\ar{d}{s} & A \\
      GE\ar[dashed]{u}{\epsilon_E}\ar[swap]{r}{G(e)} & GA\ar[shift left=0.75ex]{r}{G(f)}\ar[swap, shift right=0.75ex]{r}{G(g)}\ar{u}{\epsilon_A} & GB\ar{u}{\epsilon_B} & {} GE\ar[swap]{r}{G(e)} & GA\mathrlap{.}\ar[swap]{ur}{\epsilon_A} \\
      E\ar{u}{u}
    \end{tikzcd}
  \end{equation*}
  La flecha inducida de $GE$ a $E$ debe ser $\epsilon_E$, de nuevo por la naturalidad de $\epsilon$.
  El diagrama de la derecha conmuta por la definición de $u$ y porque $s$ es la estructura de coálgebra de $A$,
  entonces $\epsilon_AG(e)u=e$, así que, por la propiedad universal de $E$ como igualador, $\epsilon_Eu=\id_E$. Para la coasociatividad, queremos ver que el siguiente diagrama conmuta:
  \begin{equation*}
    \begin{tikzcd}
      E\ar{r}{u}\ar[swap]{d}{u} & GE\ar{d}{\delta_E} \\
      GE\ar[swap]{r}{G(u)} & G^2E.
    \end{tikzcd}
  \end{equation*}
  Para ello consideramos el siguiente diagrama:
  \begin{equation*}
    \begin{tikzcd}
      A\ar{r}{s} & GA\ar[shift left=0.75ex]{r}{G(s)}\ar[swap, shift right=0.75ex]{r}{\delta_A} & G^2A \\
      E\ar{u}{e}\ar[swap]{r}{u} & GE\ar[shift left=0.75ex]{r}{G(u)}\ar[swap, shift right=0.75ex]{r}{\delta_E}\ar{u}{G(e)} & G^2E.\ar[swap]{u}{G^2(e)}
    \end{tikzcd}
  \end{equation*}
  El cuadrado de la izquierda conmuta por la construcción de $u$ y ambos cuadrados derechos conmutan, en un caso
  porque $G$ es funtor y en el otro caso porque $\delta$ es una transformación natural. Más aún, $\delta_As=G(s)s$
  por la coasociatividad de $A$, pero $G^2(e)$ es un igualador y como los cuadros conmutan, $\delta_Ase$, $G(s)se$
  también igualan a $G^2(f)$ con $G^2(g)$, entonces por la propiedad universal
  de $G^2E$ obtenemos que $\delta_Eu=G(u)u$. Ahora, para demostrar que $(E,u)$ es el
  igualador de $f$ y $g$ en $\mathcal{E}_\textbf{G}$ nos fijamos en el siguiente diagrama.
  \begin{equation*}
    \begin{tikzcd}
      E\ar{r}{e}\ar[swap]{d}{u} & A\ar[shift left=0.75ex]{r}{f}\ar[swap, shift right=0.75ex]{r}{g}\ar{d}{s} & B\ar{d}{t} \\
      GE\ar[swap]{r}{G(e)} & GA\ar[shift left=0.75ex]{r}{G(f)}\ar[swap, shift right=0.75ex]{r}{G(g)} & GB.
    \end{tikzcd}
  \end{equation*}
  Claramente $e$, visto como morfismo de coálgebras, iguala a $f$ con $g$. Sea $h\colon (X,w)\to (A,s)$ otro morfismo
  de coálgebras que iguala a $f$ y a $g$, entonces por la propiedad universal de
  $E$ y $GE$ existe un único morfismo $j\colon X\to E$ tal que $ej=h$ y un único $k\colon GX\to GE$ tal que $G(e)k=G(h)$. Entonces tenemos al diagrama:
  \begin{equation*}
    \begin{tikzcd}
      X\ar[dashed]{r}{j}\ar[swap]{d}{w} & E\ar{d}{u} \\
      GX\ar[swap,dashed]{r}{k} & GE.
    \end{tikzcd}
  \end{equation*}
  Por construcción, $ej=h$ entonces $G(e)G(j)=G(h)$, lo cual implica que $G(j)=k$ y si el diagrama anterior conmutara, 
  $j$ sería un morfismo de coálgebras. Para demostrar que conmuta nos fijamos en el diagrama siguiente:
  \begin{equation*}
    \begin{tikzcd}
      {} & X\ar{rr}{w}\ar[swap]{ld}{h}\ar[crossing over, near end]{dd}{j} & {} & GX\ar{ld}{G(h)}\ar{dd}{G(j)} \\
      A\ar[near start]{rr}{s} & {} & GA & {} \\
      {} & E\ar{lu}{e}\ar[swap]{rr}{u} & {} & GE\ar[tail]{lu}{G(e)} 
    \end{tikzcd}
  \end{equation*}
  Ya sabemos que cada cara del prisma conmuta, salvo por la frontal. Sin embargo, como $G(e)$ es un igualador
  y por tanto mono, si seguimos el diagrama es fácil verificar que $uj=G(j)w$. Por tanto, existe un morfismo 
  $j\colon (X,w)\to (E,u)$ de coálgebras tal que $ej=h$ en $\mathcal{E}_\textbf{G}$ y es único por la propiedad universal de $E$ como igualador. 
  Con esto concluimos que $\mathcal{E}_\textbf{G}$ tiene límites finitos.
\end{proof}

\begin{lemma}
  La categoría $\mathcal{E}_\textbf{G}$ tiene exponenciales.
\end{lemma}
\begin{proof}
  Seguiremos la demostración de Mac Lane y Moerdijk en
  \cite{maclane2012sheaves}.  Recordemos que tenemos una adjunción entre el funtor olvidadizo $U\colon \mathcal{E}_\textbf{G}\to \mathcal{E}$, y el
  funtor colibre $R\colon \mathcal{E}\to \mathcal{E}_\textbf{G}$.
  La coálgebra colibre de un objeto $X$ es $(GX, \delta_X)$ y un morfismo $h\colon U(Y,\alpha)\to X$ es enviado a
  $G(h)\alpha\colon (Y,\alpha)\to RX$ bajo la adjunción. Sean $(A,s), (B,t)$ y $(C,u)$ coálgebras, entonces tenemos la siguiente cadena de isomorfismos naturales:
  
  \begin{equation*}
    \homs_{\mathcal{E}}(A\times B, C) \cong \homs_{\mathcal{E}}(U(A,s), C^B) \cong \homs_{\mathcal{E}_\textbf{G}}((A,s), R(C^B)).
  \end{equation*}
  Por el isomorfismo anterior, un buen candidato
  para el exponencial $(C,u)^{(B,t)}$ es $R(C^B)$, sin embargo, debemos restringirnos a los morfismos
  de coálgebras, por lo que sospechamos que hay un subobjeto $(E,v)$ de $R(C^B)$ tal que tengamos
  el siguiente isomorfismo natural:
  \begin{equation*}
    \homs_{\mathcal{E}_\textbf{G}}((A,s)\times (B,t), (C,u)) \cong \homs_{\mathcal{E}_\textbf{G}}((A,s), (E,v)).
  \end{equation*}
  Lo anterior implicaría que $(E,v)$ es el exponencial que buscamos. Sea $f\colon (A,s)\times (B,t) \to (C,u)$,
  definimos a $f'$ como su transpuesta exponencial en $\mathcal{E}$ y sea $f''=G(f')s$. Si traducimos
  la condición de que $f$ sea morfismo de coálgebras en una condición sobre $f''$, sabremos como
  definir a $(E,v)$. Recordando que para cualesquiera $X,Y$ en $\mathcal{E}$ hay un isomorfismo 
  $\phi\colon G(X\times Y)\to GX\times GY$ natural en $X,Y$ y utilizando la construcción del producto en el lema anterior, la condición de que $f$ sea un morfismo de 
  coálgebras se expresa con el siguiente diagrama:
  \begin{equation*}
    \begin{tikzcd}
      A\times B\ar{r}{s\times t}\ar[swap]{d}{f} & GA\times GB\ar{r}{\phi^{-1}} & G(A\times B)\ar{d}{G(f)} \\
      C\ar{rr}{u} & {} & GC.
    \end{tikzcd}
  \end{equation*}
  Ahora transponemos cada morfismo empezando por $G(f)\phi^{-1}(s\times t)$ con la siguiente observación:
  \begin{equation*}
    \begin{tikzcd}
      A\times B\ar[equals]{d}\ar{r}{s\times 1_B}\ar[phantom]{rd}{(a)} & GA\times B\ar[swap]{d}{Gf'\times 1_B}\ar{r}{1_{GA}\times t}\ar[phantom]{rd}{(b)} & GA\times GB\ar[swap]{d}{Gf'\times 1_{GB}}\ar{r}{\cong}\ar[phantom]{rd}{(c)} & G(A\times B)\ar[swap]{d}{G(f'\times 1_B)}\ar{r}{Gf}\ar[phantom]{rd}{(d)} & GC\ar[equals]{d}\\
      A\times B\ar[swap]{r}{f''\times 1_B}              & G(C^B)\times B\ar[swap]{r}{1\times t}                       & G(C^B)\times GB\ar[swap]{r}{\phi^{-1}}                 & G(C^B\times B)\ar[swap]{r}{G(ev_{B,C})}             & GC,
    \end{tikzcd}
  \end{equation*}
  donde la fila de arriba es precisamente $G(f)\phi^{-1}(s\times t)$ y cada cuadro conmuta, en (a) por la definición de $f''$, (b) es inmediato, (c) es por la 
  naturalidad de $\phi$ y (d) conmuta pues es aplicar al funtor $G$ a la definición de $f'$ como 
  transpuesta. La conmutatividad de (d), luego (c), etcétera, nos permite concluir que la fila de
  arriba es igual a la fila de abajo, que es un morfismo más sencillo de transponer
  \begin{equation*}
    \overline{G(ev)\phi^{-1}(1\times t)(f''\times 1_B)} = \overline{G(ev)\phi^{-1}(1\times t)}f'' = (GC)^t\overline{G(ev)\phi^{-1}}f''.
  \end{equation*}
  Por otro lado, $\overline{uf} = u^Bf'$. Si nombramos $p=\overline{G(ev)\phi^{-1}}$, que $f$ sea
  morfismo de coálgebras es equivalente a que $u^Bf' = (GC)^tpf''$. Ahora pasamos la igualdad a
  través de la adjunción
  \begin{align*}
    G(u^Bf')s &= G(u^B)G(f')s = G(u^B)f'', \\
    G((GC)^tpf'')s &= G((GC)^t)G(p)G(f'')s = G((GC)^t)G(p)\delta_{C^B}f''.
  \end{align*}
  Recordando que $f''\colon (A,s)\to (G(C^B),\delta_{C^B})$. Como $f''$ iguala a los anteriores
  morfismos entre coálgebras colibres, aprovechamos que $\mathcal{E}_\textbf{G}$ tiene límites finitos
  y proponemos a $(E,v)$ como el igualador de $G(u^B)$ con $G((GC)^t)G(p)\delta_{C^B}$
  \begin{equation*}
    \begin{tikzcd}
      (E,v)\ar{r}{i} & (G(C^B),\delta_{C^B})\ar{r}{G(u^B)}\ar{d}{\delta_{C^B}} & (G((GC)^B),\delta_{(GC)^B}) \\
      (A,s)\ar[dashed]{u}{f^*}\ar{ru}[description]{f''} & (G(G(C^B)),\delta_{G(C^B)})\ar{r}{G(p)} & (G((GC)^{GB}),\delta_{(GC)^{GB}}).\ar{u}{G((GC)^t)} \\
    \end{tikzcd}
  \end{equation*}
  El mapeo de $f$ a $f^*$ define un isomorfismo natural ya que cada asignación se hizo a través
  de una propiedad universal, por lo que $(E,v)\cong (C,u)^{(B,t)}$.
\end{proof}
Procedemos a mostrar que $\mathcal{E}_\textbf{G}$ tiene un clasificador de subobjetos, para lo cual necesitaremos dos lemas.
\begin{lemma}
  Sean (A,s) una coálgebra y $m\colon D\rightarrowtail A$ un subobjeto de $A$ en $\mathcal{E}$. Entonces hay
  a lo más un morfismo $d\colon D\to GD$ tal que $(D,d)$ sea coálgebra y $m$ un morfismo de coálgebras.
\end{lemma}
\begin{proof}
  Supongamos que existe un $d\colon D\to GD$ tal que $(D,d)$ sea coálgebra y $m$ un morfismo de coálgebras. Entonces tenemos
  el siguiente diagrama conmutativo
  \begin{equation}
    \label{eq:subalgebra}
    \begin{tikzcd}
      D\ar[tail]{r}{m}\ar[swap]{d}{d} & A\ar{d}{s} \\
      GD\ar[swap]{r}{G(m)} & GA.
    \end{tikzcd}
  \end{equation}
  Como $m$ es mono y $G$ es exacto por la izquierda, entonces $G(m)$ es mono. Por lo tanto, $d$ debe ser único.
\end{proof}
\begin{lemma}
  Si $(A,s)$ es coálgebra, $m\colon D\to A$ es un monomorfismo y $d\colon D\to
  GD$ cumple que $G(m)d=sm$, entonces $(D,d)$ 
  es coálgebra y el diagrama~\eqref{eq:subalgebra} es un producto fibrado en $\mathcal{E}$.
\end{lemma}
\begin{proof}
  Veamos que $(D,d)$ es coálgebra, para ello consideremos el siguiente diagrama:
  \begin{equation*}
    \begin{tikzcd}
      D\ar{rrr}{d}\ar[swap]{ddd}{d}\ar{rd}{m} & {} & {} & GD\ar{ddd}{Gd}\ar{ld}{Gm} \\
      {} & A\ar{r}{s}\ar{d}{s} & GA\ar{d}{Gs} & {} \\
      {} & GA\ar{r}{\delta_A} & G^2A & {} \\
      GD\ar{rrr}{\delta_D}\ar{ru}{Gm} & {} & {} & G^2D.\ar{lu}{G^2m}
    \end{tikzcd}
  \end{equation*}
  El cuadrado interior conmuta dado que $(A,s)$ es coálgebra, el trapecio inferior conmuta por la
  naturalidad de $\delta$ y los demás trapecios conmutan por la hipótesis de que $G(m)d=sm$. Para que
  el exterior conmute, basta ver que $G^2(m)G(d)d=G^2(m)\delta_Dd$, ya que $G^2(m)$ es mono, entonces
  tenemos lo siguiente:
  \begin{align*}
    G^2(m)G(d)d &= G(s)G(m)d = G(s)sm, \\
    G^2(m)\delta_Dd &= \delta_A G(m)d = \delta_Asm.
  \end{align*}
  Como el interior conmuta, concluimos que el exterior también y por tanto, $d$ cumple la coasociatividad.
  Para la counidad, consideremos el siguiente diagrama:
  \begin{equation*}
    \begin{tikzcd}
      D\ar{r}{d}\ar[swap,tail]{d}{m} & GD\ar{r}{\epsilon_D}\ar[tail]{d}{Gm} & D\ar[tail]{d}{m} \\
      A\ar{r}{s} & GA\ar{r}{\epsilon_A} & A.
    \end{tikzcd}
  \end{equation*}
  El cuadrado de la izquierda conmuta por hipótesis y el de la derecha por la naturalidad de $\epsilon$,
  esto basta para mostrar que el exterior conmuta. Como $(A,s)$ es coálgebra, $\epsilon_As=\id_A$, entonces 
  $m\epsilon_Dd=\id_Am=m \id_D$, pero $m$ es mono así que $\epsilon_Dd=\id_D$ y $d$ cumple la counidad.
  Observemos con más cuidado el diagrama anterior pues resulta que el cuadrado izquierdo es un producto fibrado.
  Supongamos que tenemos un objeto $X$ y morfismos $p\colon X\to GD$ y $q\colon
  X\to A$ tales que el exterior del siguiente diagrama conmuta
  \begin{equation*}
    \begin{tikzcd}
      X\ar[bend left=30]{rrd}{p}\ar[bend right=30]{rdd}[swap]{q}
          \ar[dashed]{rd}{\epsilon_Dp}\\[-2ex]
      & [-2ex] D\ar{r}{d}\ar[swap]{d}{m} & GD\ar{d}{Gm} \\
      & A\ar[swap]{r}{s} & GA.
    \end{tikzcd}
  \end{equation*}
  Notemos que:
  \begin{align*}
    G(m)d\epsilon_Dp = sm\epsilon_Dp = s\epsilon_AG(m)p &= s\epsilon_Asq = sq = G(m)p \\
    m\epsilon_Dp = \epsilon_AG(m)p &= \epsilon_Asq = q.
  \end{align*}
  Dado que $G(m)$ es mono, nuestra propuesta $\epsilon_Dp$ hace conmutar a los triángulos deseados
  y es único porque $m$ es mono. Por tanto, el diagrama es un producto fibrado.
\end{proof}
\begin{lemma}
  La categoría $\mathcal{E}_\textbf{G}$ tiene un clasificador de subobjetos.
\end{lemma}
\begin{proof}
  Empezamos con $v\colon 1\to \Omega$ el clasificador de subobjetos en $\mathcal{E}$ y definimos a $\tau\colon G\Omega\to \Omega$
  como el morfismo que clasifica al mono $G(v)\colon G1\to G\Omega$. Sean $(A,s),(D,d)$ coálgebras con $m\colon (D,d)\rightarrowtail (A,s)$
  un monomorfismo y sea $h\colon A\to \Omega$ el morfismo que clasifica a $m$ en $\mathcal{E}$. Entonces
  tenemos el siguiente diagrama:
  \begin{equation*}
    \label{eq:clasificadorm}
    \begin{tikzcd}
      D\ar{r}{d}\ar[swap]{d}{m} & GD\ar{r}{G(!_D)}\ar{d}{Gm} & G1\ar{r}{!_{G1}}\ar{d}{Gv} & 1\ar{d}{v} \\
      A\ar{r}{s} & GA\ar{r}{Gh} & G\Omega\ar{r}{\tau} & \Omega\mathrlap{,} \\
    \end{tikzcd}
  \end{equation*}
  donde no solo cada cuadro conmuta, sino que es un producto fibrado. El cuadrado de la izquierda debido
  al lema anterior, el cuadrado de en medio es $G$ aplicado a la construcción de $h$, pero $G$ preserva
  productos fibrados y el cuadrado de la derecha es producto fibrado por la definición de $\tau$.
  Por lo tanto, el exterior es un producto fibrado y $\tau G(h)s$ clasifica a $m$. Entonces $h=\tau G(h)s$,
  donde $G(h)s$ es justamente la transposición de $h$ bajo la adjunción $U\dashv R$. Más aún,
  \begin{align*}
    \id_{G\Omega}G(h)s = G(\tau G(h)s) s &= G(\tau)G^2(h)G(s)s \\
    &=G(\tau)G^2(h)\delta_A s  && \text{Coasociatividad de s} \\
    &=G(\tau)\delta_{\Omega}G(h)s && \text{Naturalidad de $\delta$}.
  \end{align*}
  Es decir, $G(h)s$ iguala a $\id_{G\Omega}$ con $G(\tau)\delta_{\Omega}$, que son morfismos de coálgebras,
  así que $G(h)s$ se factoriza a través del igualador
  \begin{equation*}
    \begin{tikzcd}
      (\Omega_G,u)\ar{r}{e} & (G\Omega,\delta_{\Omega})\ar[shift left=0.75ex]{r}{id_{G\Omega}}\ar[swap, shift right=0.75ex]{r}{G(\tau)\delta_{\Omega}} & (G\Omega,\delta_{\Omega}) \\
      (A,s).\ar[dashed]{u}{h^*}\ar[swap]{ru}{G(h)s}
    \end{tikzcd}
  \end{equation*}
  Esta factorización se da para cualquier monomorfismo de coálgebras, en particular, como $G$
  es exacto por la izquierda, hay una coálgebra trivial $(1,\psi)$ con el isomorfismo $\psi\colon 1\to G1$ y
  una sub-coálgebra $!\colon (1,\psi)\rightarrowtail (1,\psi)$ cuyo morfismo característico en $\mathcal{E}$ es $v\colon 1\to \Omega$.
  Por tanto, $G(v)\psi$ se factoriza a través de $\Omega_G$ por un morfismo, digamos, $v_G$.
  \begin{equation*}
    \begin{tikzcd}
      D\ar{r}{!_D}\ar[swap]{d}{m} & 1\ar{r}{\cong}\ar{d}{v_G} & G1\ar{d}{Gv} \\
      A\ar{r}{h^*}\ar[to path = { -- ([yshift=-3ex]\tikztostart.south) -| (\tikztotarget.south)\tikztonodes}]{rr}[swap,pos=0.25]{G(h)s}
      & \Omega_G\ar{r}{e} & G\Omega.
    \end{tikzcd}
  \end{equation*}
Ahora el cuadrado de la derecha es un producto fibrado al tener un isomorfismo arriba y un mono abajo,
mientras que el exterior es un producto fibrado al ser un rectángulo de~\eqref{eq:clasificadorm} (hay una única
flecha de $D$ a $G1\cong 1$), por tanto, el cuadrado de la izquierda es un producto fibrado. De esta manera,
para que $v_G\colon 1\to \Omega_G$ sea el clasificador de subobjetos en $\mathcal{E}_\textbf{G}$, solo falta mostrar que
$h^*$ es único. Supongamos que hay otro morfismo $f\colon (A,s)\to (\Omega_G,u)$ tal que el diagrama de la izquierda
es un producto fibrado. Para mostrar que $f=h^*$, basta mostrar que $ef=eh^*=G(h)s$, pero $G(h)s$ es un morfismo
hacia una coálgebra colibre, por lo que, pasando por la adjunción $U\dashv R$, es suficiente que
$\epsilon_{\Omega}ef=h$ en $\mathcal{E}$.
\begin{align*}
  \epsilon_{\Omega}ef &= \epsilon_{\Omega}G(\tau)\delta_{\Omega}ef && \text{Definición de e} \\
  &= \tau\epsilon_{G\Omega}\delta_{\Omega}ef && \text{Naturalidad de $\epsilon$} \\
  &= \tau ef. && \text{Estructura de comónada}
\end{align*}
Sin embargo tenemos el siguiente diagrama:
\begin{equation*}
  \begin{tikzcd}
    D\ar{r}{!_D}\ar[swap]{d}{m} & 1\ar{r}{\cong}\ar{d}{v_G} & G1\ar{r}{!_{G1}}\ar{d}{Gv} & 1\ar{d}{v} \\
    A\ar{r}{f} & \Omega_G\ar{r}{e} & G\Omega\ar{r}{\tau} & \Omega.
  \end{tikzcd}
\end{equation*}
Dado que $f$ clasifica a $m$, el cuadrado de la izquierda es un producto fibrado en $\mathcal{E}_\textbf{G}$, pero por
como construimos los límites finitos en un lema anterior, el funtor que olvida los preserva, entonces también es un
producto fibrado en $\mathcal{E}$. El cuadro de en medio tiene un iso arriba y un mono abajo, mientras que el cuadro
de la derecha es un producto fibrado al ser la definición de $\tau$. Lo anterior nos permite concluir que el exterior
es un producto fibrado y que $\tau ef$ clasifica a $m$ en $\mathcal{E}$, en otras palabras, $\tau ef = h$ que era
lo que necesitábamos para que $f=h^*$. Por lo tanto, $\mathcal{E}_G$ tiene clasificador de subobjetos y es un topos.
\end{proof}











